\section{Multi-Agent Systems and Distributed Inference}
\label{sec:multi-agent}

If each agent is performing BP over its local factor graph, then a
multi-agent system can be interpreted as a distributed inference computation
over a larger global factor graph. This resembles loopy belief propagation
on a graph too large to fit in a single model.

This is an interpretive framework, not a proven equivalence. Real LLM
agents are not calibrated BP nodes. Communication between agents is lossy
text, not exact belief messages. The convergence guarantees from formal
loopy BP do not transfer directly to multi-agent LLM systems.

\subsection*{What the BP Lens Suggests}

Even without exact equivalence, the distributed BP framing suggests useful
design principles: agent boundaries may work better when aligned with factor
graph structure; message formats carrying structured uncertainty may
coordinate better than raw text; convergence criteria inspired by BP belief
stability could inform stopping rules for multi-agent deliberation.

\subsection*{What Would Be Required for a Formal Account}

For the distributed BP interpretation to become a formal account, agents
would need to maintain calibrated beliefs over a shared declared factor
graph, pass approximately correct belief messages, update beliefs according
to BP rules, and operate over a graph structure where loopy BP converges.
None of these conditions hold for current LLM multi-agent systems.

\subsection*{Epistemic Status}

\begin{center}
\begin{tabular}{ll}
\toprule
Claim & Status \\
\midrule
Multi-agent systems resemble distributed BP & Mechanistic interpretation \\
BP lens suggests useful design principles & Architectural proposal \\
Current LLM agents converge to correct posteriors & Not supported \\
BP-inspired coordination protocols may help & Speculative hypothesis \\
\bottomrule
\end{tabular}
\end{center}